El intervalo exacto de posibilidades es
\[
e^{-3/2}-e^{-2}\le x(2)\le e^{-1}-e^{-3/2}.
\]
El extremo inferior se alcanza aportando potencia uno en \([0,1/2]\); el superior, en \([1/2,1]\), con potencia cero fuera del tramo elegido. La temperatura retiene más del aporte tardío porque ha tenido menos tiempo para perderlo.
